% -----------------------------------------------------------
\section{Condition}
% -----------------------------------------------------------
	\label{CONDITION}
	
% % ----------------------
% \subsection{See also}
% % ----------------------

% ----------------------
\subsection{1828 American Dictionary of the English Language} % *1828
% ----------------------
	\label{CONDITION-1828}
	
	% -----
	\subsubsection{Noun}
	% -----
	
		The first five meanings here give a meaning of ``condition'' where it is synonymous with ``situation'' or ``circumstance.'' The other three meanings have a more legal flavor and are much more interesting for the scriptures.

		\begin{compactitem}
			\item[1] State; a particular mode of being; applied to external circumstances, to the body, to the mind, and to things.
			\item[2] Quality; property; attribute.
			\item[3] State of mind; temper; temperament; complexion.
			\item[4] Moral quality; virtue or vice.
			\item[5] Rank, that is, state with respect to the orders or grades of society, or to property
			\item[6] Terms of a contract or covenant; stipulation; that is, that which is set, fixed, established or proposed.
			\item[7] A clause in a bond, or other contract containing terms or a stipulation that it is to be performed, and in case of failure, the penalty of the bond is to be incurred.
			\item[8] Terms given, or provided, as the ground of something else; that which is established, or to be done, or to happen, as requisite to another act
		\end{compactitem}
		
	% -----
	\subsubsection{Intransitive verb}
	% -----
	
		\begin{compactitem}
			\item To make terms; to stipulate.
		\end{compactitem}
		
	% -----
	\subsubsection{Transitive verb}
	% -----
	
		\begin{compactitem}
			\item To contract; to stipulate.
		\end{compactitem}

% % ----------------------
% \subsection{Oxford English Dictionary} % *OED
% % ----------------------
	% \label{CONDITION-OED}

	% \begin{compactitem}
		% \item[]
	% \end{compactitem}

% ----------------------
\subsection{Scriptural Usage} % *SCRIPTURES
% ----------------------
	\label{CONDITION-SCRIPTURES}

	% -----
	\subsubsection{Old Testament} % *OT
	% -----
		\label{CONDITION-OT}
	
		% --
		\paragraph{KJV}
		% --
			\label{CONDITION-OT-KJV}
	
			\begin{tabular}{ll}
				Total occurrences in the OT & 1
			\end{tabular}
			
			\begin{compactdesc}
				\item[1 Samuel 11:2] And Nahash the Ammonite answered them, On this condition will I make a covenant with you, that I may thrust out all your right eyes, and lay it for a reproach upon all Israel. \textit{(``Condition'' is in italics here, and so are ``a covenant'' and ``for'' at the end.)}
			\end{compactdesc}
			
		% --
		\paragraph{Hebrew Bible}
		% --
			\label{CONDITION-HEBREW}
		
			The word ``condition'' in 1 Samuel 11:2 does not represent a Hebrew word, and is added just to make the passage more comprehensible in English. The Hebrew just says \he{b*:zo't}, ``by this'' or ``with this.''
			
		% --
		\paragraph{Septuagint}
		% --
			\label{CONDITION-LXX}
		
			In the part corresponding to ``condition'' in the KJV, 1 Samuel 11:2 in the LXX has \gr{ἐν ταύτῃ}.
		
	% -----
	\subsubsection{New Testament} % *NT
	% -----
		\label{CONDITION-NT}
	
		% --
		\paragraph{KJV}
		% --
			\label{CONDITION-NT-KJV}
			
	
			\begin{tabular}{ll}
				Total occurrences in the NT & 1
			\end{tabular}
			
			\begin{compactdesc}
				\item[Luke 14:32] Or else, while the other is yet a great way off, he sendeth an ambassage, and desireth conditions of peace.
			\end{compactdesc}
			
		% --
		\paragraph{Greek New Testament} % *GREEK
		% --
			\label{CONDITION-GREEK}
		
			Similarly to the OT in Hebrew and Greek, the NT Greek uses ``conditions'' to flesh out the meaning of Luke 14:32. The Greek just has \gr{τὰ πρὸς εἰρήνην}.
		
	% -----
	\subsubsection{Book of Mormon} % *BOM
	% -----
		\label{CONDITION-BOM}
	
		\begin{tabular}{ll}
			Total occurrences in the BOM & 14
		\end{tabular}
		
		\begin{compactdesc}
			\item[Mosiah 4:8] And this is the means whereby salvation cometh.  And there is none other salvation save this which hath been spoken of; neither are there any conditions whereby man can be saved except the conditions which I have told you.
			\item[Mosiah 19:15] Therefore the Lamanites did spare their lives, and took them captives and carried them back to the land of Nephi, and granted unto them that they might possess the land, under the conditions that they would deliver up king Noah into the hands of the Lamanites, and deliver up their property, even one half of all they possessed, one half of their gold, and their silver, and all their precious things, and thus they should pay tribute to the king of the Lamanites from year to year.
			\item[Alma 5:10] And now I ask of you on what conditions are they saved?  Yea, what grounds had they to hope for salvation?  What is the cause of their being loosed from the bands of death, yea, and also the chains of hell?
			\item[Alma 17:15] Thus they were a very indolent people, many of whom did worship idols, and the curse of God had fallen upon them because of the traditions of their fathers; notwithstanding the promises of the Lord were extended unto them on the conditions of repentance.
			\item[Alma 27:24] And now behold, this will we do unto our brethren, that they may inherit the land Jershon; and we will guard them from their enemies with our armies, on condition that they will give us a portion of their substance to assist us that we may maintain our armies.
			\item[Alma 42:13] Therefore, according to justice, the plan of redemption could not be brought about, only on conditions of repentance of men in this probationary state, yea, this preparatory state; for except it were for these conditions, mercy could not take effect except it should destroy the work of justice.  Now the work of justice could not be destroyed; if so, God would cease to be God.
			\item[Alma 44:11] Now I cannot recall the words which I have spoken, therefore as the Lord liveth, ye shall not depart except ye depart with an oath that ye will not return again against us to war.  Now as ye are in our hands we will spill your blood upon the ground, or ye shall submit to the conditions which I have proposed.
			\item[Alma 54:11] But behold, it supposeth me that I talk to you concerning these things in vain; or it supposeth me that thou art a child of hell; therefore I will close my epistle by telling you that I will not exchange prisoners, save it be on conditions that ye will deliver up a man and his wife and his children, for one prisoner; if this be the case that ye will do it, I will exchange.
			\item[Helaman 5:11] And he hath power given unto him from the Father to redeem them from their sins because of repentance; therefore he hath sent his angels to declare the tidings of the conditions of repentance, which bringeth unto the power of the Redeemer, unto the salvation of their souls.
			\item[Helaman 14:11] And ye shall hear my words, for, for this intent have I come up upon the walls of this city, that ye might hear and know of the judgments of God which do await you because of your iniquities, and also that ye might know the conditions of repentance;
			\item[Helaman 14:18] Yea, and it bringeth to pass the condition of repentance, that whosoever repenteth the same is not hewn down and cast into the fire; but whosoever repenteth not is hewn down and cast into the fire; and there cometh upon them again a spiritual death, yea, a second death, for they are cut off again as to things pertaining to righteousness.
			\item[Helaman 16:11] And these were the conditions also, in the eighty and eighth year of the reign of the judges.
		\end{compactdesc}
		
	% -----
	\subsubsection{Doctrine and Covenants} % *DC
	% -----
		\label{CONDITION-DC}
	
		\begin{tabular}{ll}
			Total occurrences in the DC & 11
		\end{tabular}
		
		\begin{compactdesc}
			\item[DC 18:12] And he hath risen again from the dead, that he might bring all men unto him, on conditions of repentance.
			\item[DC 46:15] And again, to some it is given by the Holy Ghost to know the differences of administration, as it will be pleasing unto the same Lord, according as the Lord will, suiting his mercies according to the conditions of the children of men.
			\item[DC 88:38--39] And unto every kingdom is given a law; and unto every law there are certain bounds also and conditions. 39 All beings who abide not in those conditions are not justified.
			\item[DC 109:65] And cause that the remnants of Jacob, who have been cursed and smitten because of their transgression, be converted from their wild and savage condition to the fulness of the everlasting gospel;
			\item[DC 113:10] We are to understand that the scattered remnants are exhorted to return to the Lord from whence they have fallen; which if they do, the promise of the Lord is that he will speak to them, or give them revelation.  See the 6th, 7th, and 8th verses.  The bands of her neck are the curses of God upon her, or the remnants of Israel in their scattered condition among the Gentiles.
			\item[DC 132:5] For all who will have a blessing at my hands shall abide the law which was appointed for that blessing, and the conditions thereof, as were instituted from before the foundation of the world.
			\item[DC 132:7] And verily I say unto you, that the conditions of this law are these: All covenants, contracts, bonds, obligations, oaths, vows, performances, connections, associations, or expectations, that are not made and entered into and sealed by the Holy Spirit of promise, of him who is anointed, both as well for time and for all eternity, and that too most holy, by revelation and commandment through the medium of mine anointed, whom I have appointed on the earth to hold this power (and I have appointed unto my servant Joseph to hold this power in the last days, and there is never but one on the earth at a time on whom this power and the keys of this priesthood are conferred), are of no efficacy, virtue, or force in and after the resurrection from the dead; for all contracts that are not made unto this end have an end when men are dead.
			\item[DC 132:17] For these angels did not abide my law; therefore, they cannot be enlarged, but remain separately and singly, without exaltation, in their saved condition, to all eternity; and from henceforth are not gods, but are angels of God forever and ever.
			\item[DC 138:19] And there he preached to them the everlasting gospel, the doctrine of the resurrection and the redemption of mankind from the fall, and from individual sins on conditions of repentance.
			\item[DC OD1:16] 
		\end{compactdesc}
		
	% -----
	\subsubsection{Pearl of Great Price} % *PGP
	% -----
		\label{CONDITION-PGP}
	
		\begin{tabular}{ll}
			Total occurrences in the PGP & 0
		\end{tabular}
		
% ----------------------
\subsection{Key Phrases} % *KEYPHRASES *PHRASES
% ----------------------
	\label{CONDITION-PHRASES}

	\begin{compactdesc}
		\item[condition(s) of repentance] \textbf{7} | \hyperref[ALMA-17-15]{Alma 17:15}; \hyperref[ALMA-42-13]{42:13}; \hyperref[HELAMAN-5-11]{Helaman 5:11}; \hyperref[HELAMAN-14-11]{14:11}, \hyperref[HELAMAN-14-18]{14:18}; \hyperref[DC18-12]{DC 18:12}; \hyperref[DC138-19]{DC 138:19}
			\begin{compactdesc}
				\item[Commentary] See \hyperref[CONDITION-repentance]{this study} below, and consider this idea in light of \hyperref[DC88-38]{DC 88:38}.
			\end{compactdesc}
	\end{compactdesc}
	
% % ----------------------
% \subsection{Bible Dictionaries} % *DICTIONARIES
% % ----------------------
	% \label{CONDITION-BD}

% % ----------------------
% \subsection{Restoration Usage} % *RESTORATION
% % ----------------------
	% \label{CONDITION-RESTORATION}

	% % -----
	% \subsubsection{Joseph Smith} % *JS
	% % -----
		% \label{CONDITION-JS}
	
		% \begin{compactdesc}
			% \item[DATE/SOURCE]
		% \end{compactdesc}
	
	% % -----
	% \subsubsection{Brigham Young} % *BY
	% % -----
		% \label{CONDITION-BY}
	
		% \begin{compactdesc}
			% \item[DATE/SOURCE]
		% \end{compactdesc}
	
% ----------------------
\subsection{Commentary and Studies} % *COMMENTARY *STUDIES
% ----------------------
	\label{CONDITION-COMMENTARY}

	% -----
	\subsubsection{Key Points}
	% -----
		\label{CONDITION-KEYS}
	
		\begin{compactitem}
			\item The ``conditions of repentance'' are the terms on which repentance works. The scriptures mention several of these conditions---see \hyperref[CONDITION-repentance]{below} for citations:
				\begin{compactitem}
					\item Whoever repents ``is not hewn down and cast into the fire.''
					\item The promises of the Lord are extended to us when we repent.
					\item The plan of redemption works only through repentance.
					\item Mercy cannot take effect in our lives without repentance.
					\item Repentance brings us to the power of the Redeemer.
					\item Individual sins are resolved through repentance.
					\item 
				\end{compactitem}
		\end{compactitem}
	
	% -----
	\subsubsection{Conditions of repentance}
	% -----
		\label{CONDITION-repentance}
		
		The word ``condition'' is part of a very important phrase in the scriptures, the ``conditions of repentance.'' This phrase occurs seven times in the scriptures, all in the BOM and DC. Here are a couple examples, first from Helaman 5:11:
		
			\begin{quote}
				And he hath power given unto him from the Father to redeem them from their sins because of repentance. Therefore he hath sent his angels to declare the tidings of the conditions of repentance, which bringeth unto the power of the Redeemer, unto the salvation of their souls.
			\end{quote}
			
		The next is Helaman 14:17--18:
		
			\begin{quote}
				But behold, the resurrection of Christ redeemeth mankind, yea, even all mankind, and bringeth them back into the presence of the Lord. 18 Yea, and it bringeth to pass the conditions of repentance, that whosoever repenteth, the same is not hewn down and cast into the fire. But whosoever repenteth not is hewn down and cast into the fire. And there cometh upon them again a spiritual death, yea, a second death, for they are cut off again as to things pertaining to righteousness. 
			\end{quote}
			
		I think the word ``condition'' in this phrase is used in its more legal sense, corresponding to meanings \#6, \#7, and \#8 \hyperref[CONDITION-1828]{above}. All three of those definitions revolve around the idea of terms in a contract. So the conditions of repentance are the terms on which the repentance ``agreement'' or law works. Helaman 14:18 actually gives a statement of these terms, or at least one of them---it is that whoever repents ``is not hewn down and cast into the fire.'' That's a condition, or term in the contract, according to which repentance works.

		There are other conditions as well. \hyperref[ALMA-17-15]{Alma 17:15} gives one, that the promises of the Lord are extended to us when we repent. \hyperref[ALMA-42-13]{Alma 42:13} gives us two more: that the plan of redemption works only through repentance, and mercy cannot take effect in our lives without repentance. Helaman 5:11 adds that repentance brings us to the power of the Redeemer; \hyperref[DC138-19]{DC 138:19} says that individual sins are resolved through repentance.
		
		So the conditions are the terms according to which repentance works. Repentance is a rule- or law-governed activity, and the conditions are those rules or laws which govern it. When we know the rules, we can receive the good that comes from repentance (meaning \#8 \hyperref[CONDITION-1828]{above}). When we know the rules, we know how to repent.
		
		We could probably add other conditions to the list as well. The frequent instruction to be baptized for repentance probably fits in here, especially as stated to clearly and forcefully in a place like \hyperref[3NEPHI-27-20]{3 Nephi 27:20}. See also the \hyperref[REMISSION-commandments]{study} on remission of sin and fulfilling the commandments, especially \hyperref[REMISSION-commandments-which]{which commandments}.
		
		There's also an interesting connection to be made from Helaman 5:11 and the idea of Christ receiving power to redeem from the Father. Christ receives this power and is able to do what he does. \hyperref[MORONI-7-27]{Moroni 7:27} says that Christ claims ``of the Father his rights of mercy.'' ``Rights'' is another legal-sounding term here. See the \hyperref[RIGHT-NOUN-1828]{meanings} of that word in 1828.
	
	% -----
	\subsubsection{Conditions and laws}
	% -----
		\label{CONDITION-laws}
		
		Specifically thinking of \hyperref[DC88-38]{DC 88:38}, you can combine the two groups of meanings of \hyperref[CONDITION-1828]{condition} to understand that verse as referring to the specific situations in which the laws apply, or as giving the circumstances for application of the laws. 
		
		This topic deserves more study, however.
		
		See also \hyperref[DC132-5]{DC 132:5} and \hyperref[DC132-7]{7}, plus other verses in that section;  